\scanpage{93}
Complexes cubiques généralisés -- (Complexes \uline{triadiques} ou \uline{trialitaires} ?)

$S$ graphe strict, $A\subset\unsure{\mathcal P|\mathfrak P}_2(S)$ l'ens des arêtes. On \deleted{\illegible{1}} appelle \uline{triangle} de $S$ tout $\underline t\in\mathcal P_3(A)$ tel que $x,y\in\underline t$, $x\ne y$ implique que $x,y$ sont ``liés'' (i.e. $\{x,y\}\in A$). Soit $T$ l'ens des triangles de $S$

\uline{Axiome} 1 \quad $\forall\,a\in A,\quad\exists!\,\underline t\in T$ tel que $a\subset\underline t$

\uline{Axiome} 2 \quad $\forall\,\underline t\in T$, $\deleted{\unsure{a}}s\in S\setminus\underline t$, $s$ est lié à un élément de $t$ (unique d'après l'axiome 1\unsure{)|,|).}

\uline{NB} La conjonction de 1 et 2 s'énonce aussi
\[
\uline{\text{Axiome (C)}}\left\{\begin{aligned}
&\text{a) toute arête est contenue dans un triangle}\\
&\text{b) }\forall\,\underline t\in T,\ s\in S\setminus\underline t,\quad\exists!\,s'\in t\ \text{tel que }s\text{ lié}\\
&\quad \text{à }s'\ (\text{i.e. }\{s,s'\}\in A).
\end{aligned}\right.
\]

Pour l'instant on suppose seulement \textcircled{1} \uline{l'axiome 1 vérifié}. Supposons $A\ne\varnothing$ (sinon l'axiome 1 est vide !) donc (grâce à l'axiome 1) $T\ne\varnothing$. Soit $t_0\in T$, \inserted{on posera parfois} $t=\{a,b,c\}$, avec $a,b,c\in \deleted{\unsure{B}}S$.

On supposera de plus pour l'instant que \textcircled{2} \uline{l'axiome 2 est vérifié pour ce $t_0$}, i.e. tt élément de $S\setminus t_0$ est lié à un des él. $a,b,c$\inserted{(unique par l'axiome 1)}. Si $a\in t$, on désigne par $\widetilde E_a$ l'ens de $S\setminus t_0$ liés à $a$. Donc si $t=\{a,b,c\}$ on a bij. can.
\[
S\setminus t_0=\deleted{\unsure{\widetilde E_\alpha}}\coprod_{\alpha\in t_0}\widetilde E_\alpha
\simeq\overline E_a\amalg\overline E_b\amalg\overline E_c
\tag*{\deleted{\unsure{(1)}}}
\]
donc bij can
\[
S\simeq t_0\amalg\coprod_{\alpha\in t_0}\widetilde E_\alpha
\tag*{(1)}
\]

\deleted{D'autre part,}
% archive sheet number: 92
